\bookchapter{الفصل 51: تطوير الواجهات — مدخل إلى Jetpack Compose}

\chapterintro{الواجهة هي كل ما يراه المستخدم. قبل ظهور Compose كانت واجهات Android تُكتب بـ XML — كأنها ملء الفراغات، أمر شاق؛ وفي 2021 أطلقت جوجل Jetpack Compose: “رسم” الواجهة مباشرة بكود Kotlin، بحيث ترى النتيجة فورًا وتتحدث المعاينة لحظة تعديل سطر واحد. بحلول 2026 أصبح Compose هو الخيار السائد بلا منازع للمشاريع الجديدة. في هذا الفصل سنتعلم “الرسم”.}

\booksection{أولًا: فكرة Compose: الواجهة دالة}

فكرة Compose الجوهرية في جملة واحدة: الواجهة = دالة. كل ما يمكن عرضه على الشاشة (نص، زر، صورة) هو “دالة قابلة للتركيب” (Composable) مميزة بـ @Composable. وعند تداخلها كقطع مكعبات تتكون الواجهة.

\begin{lstlisting}[language=Kotlin]
@Composable
fun Greeting(name: String) {
    Text(text = "مرحبًا، $name!")     // "قطعة نص"
}

// تركيب متداخل: قطع كبيرة تحتوي قطعًا صغيرة
@Composable
fun HomeScreen() {
    Column {                          // مرتبة رأسيًا
        Greeting(name = "أحمد")
        Greeting(name = "سارة")
        Button(onClick = { /* ماذا يحدث عند الضغط */ }) {
            Text("اضغط هنا")
        }
    }
}
\end{lstlisting}

انتبه إلى “القبعة” @Composable — إنها تخبر المترجم “هذه قطعة من مكعبات الواجهة”. يمكن تداخل القطع إلى ما لا نهاية، وهذا كل سر Compose.

\booksection{ثانيًا: المكونات الأساسية: نص، زر، حقل إدخال، صورة}

مكونات Compose هي “قطع المكعبات”، وأكثرها شيوعًا أربعة:

\begin{center}
\begin{tabularx}{\linewidth}{|X|X|X|}
\hline
\textbf{المكون} & \textbf{ماذا يفعل} & \textbf{المعاملات الأساسية} \tabularnewline
\hline
Text & عرض نص & text (النص المراد عرضه) \tabularnewline
Button & زر & onClick (دالة الاستجابة للنقر) \tabularnewline
TextField & حقل إدخال & value و onValueChange (دالة الاستجابة لتغير الإدخال) \tabularnewline
Image & عرض صورة & painter (مورد الصورة) \tabularnewline
\hline
\end{tabularx}
\end{center}

\begin{center}\small\color{gray}المكونات الأساسية الأربعة في Compose\end{center}

\begin{lstlisting}[language=Kotlin]
@Composable
fun BasicWidgets() {
    Column {
        Text(text = "هذا نص")
        Button(onClick = { println("تم الضغط!") }) {
            Text("زر")
        }
        var input by remember { mutableStateOf("") }
        TextField(
            value = input,
            onValueChange = { input = it },
            label = { Text("اكتب شيئًا") }
        )
    }
}
\end{lstlisting}

في سطري حقل الإدخال ظهرت العبارة remember \{ mutableStateOf(“”) \} — وهذه هي “سحر الحالة” في Compose وسنفصله في القسم الخامس من هذا الفصل. اكتفِ الآن بالتعرف على شكلها.

\booksection{ثالثًا: ثلاثية التخطيط: Column و Row و Box}

الآن لدينا المكونات؛ كيف نرتبها؟ احفظ ثلاثة “صناديق”:

\begin{itemize}
  \item Column: ترتيب رأسي (مثل الوقوف في طابور)؛
  \item Row: ترتيب أفقي (مثل الجلوس جنبًا إلى جنب)؛
  \item Box: التراكب (مثل كومة أوراق اللعب، مناسب للشارات والطبقات).
\end{itemize}

\begin{lstlisting}[language=Kotlin]
Column {
    Text("السطر الأول")
    Row {
        Text("يسار")
        Text("يمين")
    }
    Box {
        Text("الطبقة السفلية")
        Text("شارة علوية متراكبة")
    }
}
\end{lstlisting}

تريد تخطيطًا أجمل؟ أضف إلى المكونات Modifier (المُزخرف) — فهو بمثابة “مكياج” للمكعبات:

\begin{lstlisting}[language=Kotlin]
Text(
    text = "نص جميل",
    modifier = Modifier
        .padding(16.dp)          // هامش من كل الجهات
        .fillMaxWidth()          // يملأ العرض كاملًا
        .background(Color.Blue)  // لون الخلفية
        .clickable { /* عند الضغط */ }// قابل للنقر
)
// استدعاءات متسلسلة: واحدًا بعد الآخر، والترتيب مهم (padding أولًا ثم الخلفية؛ الخلفية لا تشمل الهامش)
\end{lstlisting}

\begin{tcbwarn}{تنبيه: ترتيب Modifier له شأن}
padding(16.dp).background(Blue) و background(Blue).padding(16.dp) تعطيان نتيجتين مختلفتين — الأولى “فراغ ثم تلوين” (الخلفية لا تشمل الهامش)، والثانية “تلوين ثم فراغ” (الخلفية تشمل الهامش). هذا أكثر موضع يخطئ فيه المبتدئ فيرسم “حواف بيضاء غريبة”. عند الشك بدّل الترتيب وستخبرك المعاينة بالجواب فورًا.

\end{tcbwarn}

\booksection{رابعًا: الحالة: نجعل الواجهة “حية”}

الواجهة الثابتة مملة. لجعل الرقم يزيد +1 عند نقر الزر أو حقل الإدخال يقبل الكتابة، نحتاج إلى الحالة (State). قاعدة Compose: الواجهة تعكس الحالة بأمانة دائمًا، ومتى تغيرت الحالة أُعيد رسم الواجهة تلقائيًا.

\begin{lstlisting}[language=Kotlin]
@Composable
fun Counter() {
    // ‏remember + mutableStateOf: يتذكر هذه القيمة المتغيرة
    var count by remember { mutableStateOf(0) }

    Column {
        Text(text = "ضغطت $count مرات")
        Button(onClick = { count++ }) {   // يغيّر الحالة
            Text("اضغط +1")                // تتحدث الواجهة تلقائيًا
        }
    }
}
\end{lstlisting}

ثلاثة أسطر تكفي لعمل عدّاد حي. وهذه هي متعة Compose مقارنة بالأسلوب القديم: لا حاجة إلى “تحديث الواجهة” يدويًا؛ يكفي تغيير البيانات. احفظ دالتين رئيسيتين: remember (تحفظ القيمة من المسح عند إعادة تركيب الواجهة) و mutableStateOf (تخبر Compose “إذا تغيرت هذه القيمة أعيد الرسم”).

المستخدم يضغط

الحالة تتغير (count++)

Compose يلاحظ التغير

يعيد رسم الواجهة (إعادة تركيب)

\begin{center}\small\color{gray}الشكل 3: تدفق البيانات أحادي الاتجاه “الحالة الواجهة” في Compose\end{center}

\booksection{خامسًا: المظهر والوضع الداكن: إعداد واحد يسري على الكل}

ما تراه في قالب المشروع MyFirstAppTheme \{ ... \} هو المظهر (Theme). وهو مبني على مواصفات التصميم Material 3، ويساعدك على إدارة الألوان والخطوط والأشكال بشكل موحد. دعم الوضع الداكن يحتاج ثلاث خطوات فقط:

\begin{enumerate}
  \item عرّف في المظهر مجموعتي ألوان: فاتحة وداكنة (القالب يولّد ColorScheme تلقائيًا)؛
  \item عندما ينتقل النظام إلى الوضع الداكن تكتشف الدالة isSystemInDarkTheme() ذلك تلقائيًا؛
  \item استخدم ألوان المظهر المعرّفة (مثل MaterialTheme.colorScheme.primary) بدلًا من تثبيت قيم الألوان يدويًا.
\end{enumerate}

\begin{tcbinfo}{نصيحة: الألوان الديناميكية في Material 3}
ميزة “Material You” ابتداءً من Android 12 تتيح للتطبيق استخراج الألوان من خلفية جهازك. قوالب Compose تدعمها افتراضيًا (dynamicColor). سيتغير “لباس” تطبيقك تلقائيًا على هواتف المستخدمين المختلفة — هذه هي رومانسية Android الحديثة.

\end{tcbinfo}

\booksection{سادسًا: المعاينة والحركة والتحرير الحي: “القوى الخارقة” في Compose}

ثلاث “قوى خارقة” في تجربة تطوير Compose؛ متى تستخدمها لن تعود إلى الوراء:

\begin{itemize}
  \item معاينة @Preview: أضف التعليق @Preview إلى الدالة القابلة للتركيب فتعرض Android Studio الواجهة فورًا على الجانب — تعديل مباشر يريك النتيجة دون تثبيت على الهاتف في كل مرة؛
  \item التحرير الحي (Live Edit): يدعم Android Studio في 2026 "التحديث الساخن" لبعض التعديلات على المحاكي/الجهاز الحقيقي؛ تعديلات صغيرة مثل تغيير لون النص تظهر على الجهاز في ثوانٍ؛
  \item واجهات الحركة (Animation API): بضعة أسطر مثل animateFloatAsState و AnimatedVisibility تكفي لحركة سلسة للعناصر — "الإحساس الاحترافي" غالبًا يكمن في هذه العشرات من الأسطر.
\end{itemize}

\begin{lstlisting}[language=Kotlin]
@Composable
fun AnimatedBox() {
    var big by remember { mutableStateOf(false) }
    // عند تغير القيمة ينتقل بسلاسة إلى القيمة الهدف
    val size by animateDpAsState(if (big) 120.dp else 60.dp)

    Box(Modifier
        .size(size)
        .background(MaterialTheme.colorScheme.primary)
        .clickable { big = !big }
    )
}
\end{lstlisting}

\begin{tcbinfo}{نصيحة: الوظيفة أولًا ثم الحركة}
خطأ المبتدئين الشائع: “الواجهة لم تُصح بعد ومع ذلك يبدؤون بالانشغال بالحركات”. اجعل الوظيفة تعمل أولًا ثم أضف الحركات دفعة واحدة في النهاية — وهذا أيضًا مبدأ “البساطة أولًا”: الحركة زينة تُضاف، والوظيفة حاجة أساسية لا غنى عنها.

\end{tcbinfo}

\begin{tcbpractice}{لنتجرب: اصنع "زر المزاج"}
حوّل مشروع القالب إلى مبدّل مزاج: زر واحد، عند نقره يتنقل بين سعيد حزين غاضب، مع تغير النص ولون الخلفية تبعًا لذلك. تلميح: استخدم when لاختيار النص واللون حسب الحالة. عند الانتهاء أرسله لأصدقائك لتجربوه وستحصد أول موجة “واو، صرت تصنع تطبيقات!”.

\end{tcbpractice}

\begin{tcbsummary}{خلاصة الفصل}
\begin{itemize}
  \item فكرة Compose: الواجهة = دوال @Composable متداخلة كقطع مكعبات
  \item المكونات الأربعة: Text / Button / TextField / Image
  \item ثلاثية التخطيط: Column رأسي، Row أفقي، Box متراكب
  \item سلسلة Modifier لل"تزيين": padding و fillMaxWidth و background و clickable — والترتيب له معنى
  \item الحالة هي الروح: remember + mutableStateOf؛ غيّر البيانات فتُعاد رسم الواجهة تلقائيًا
  \item المظهر عبر Material 3 مع الألوان الديناميكية؛ إعداد واحد للوضع الداكن يسري على الكل
\end{itemize}

\end{tcbsummary}
